\section{Introduction}
The course task asks for EMA38 and EMA100 over five-minute tumbling windows,
followed by buy and sell advisories at crossovers, with publication restricted
to subscribed symbols~\cite{brief}. The input is a recorded week of market
events, rather than a live exchange connection~\cite{dataset}.
\pending{State the final scope, motivating question and contributions. Cite
implementation PRs and retained experiments for every headline result.}

\section{Dataset analysis}
The source covers 8--14 November 2021 and includes price, quote and
housekeeping events~\cite{dataset,brief}. The report must distinguish all
input rows from valid price events and calculated windows. Dataset analysis
has its own space because the academic-writing rubric calls for a strong
focus on event distribution~\cite{rubric}.
\pending{Select final figures from \texttt{analysis/out/} and explain their
provenance via \texttt{docs/eda-week.md} and the daily manifests. Include event
taxonomy, per-exchange rates and peaks, long-tail symbol skew, timestamp
quality and empty-window frequency. Report the selected file checksums and
subset, rather than quoting rounded brief totals as measured totals.}
\pending{Explain which measured properties justify watermark bounds and
partitioning. Attribute the dataset to Frischbier et al. and retain its
CC BY-NC-SA 4.0 notice when adapting data figures~\cite{dataset}.}

\section{System architecture}
\pending{Draw and verify the final data path against the submitted revision:
CSV replayer, Kafka ticks and subscriptions, Flink keyed EMA/crossover state,
subscription filtering, result topics, dashboard and monitoring. Label which
components share a host. The roadmap is a plan, not proof of deployment.}
\pending{Explain keying by symbol, partition load, normalization, metadata
handling, event time, watermarks, idleness and late-event policy. Explain
finite-input final-window completion and the distinction from a live stream.
Link each design trade-off to its D-number in the existing
\path{docs/decisions.md}; do not create a second architecture decision log.}
\pending{Compare the chosen state/windowing approach with measured alternatives
where available. Mark proposed ablations as unperformed until evidence exists.}

\section{Implementation}
\pending{Verify the implementation description from source and PR evidence:
price versus housekeeping classification, malformed timestamp handling,
latest-event-time close and equal-time ties, EMA initialization, empty windows,
overnight state, crossover inequalities, subscription bootstrap and updates.}
\pending{Record hardware, runtime versions, programming and GUI tools, effective
resource limits and deployment topology from each experiment. Describe restart
and replay guarantees only to the extent tested; fresh-topic replay does not
prove checkpoint recovery or exactly-once sinks.}
\pending{Insert source and PR links for the final code revision. Reuse the
interfaces and decisions documentation; validate it against code before
turning these prompts into factual prose.}

\section{Evaluation}
\subsection{Method and provenance}
\pending{Define correctness, throughput denominators and latency boundaries.
Use one completed \path{report/experiments/template.json} record per run,
with raw paths, configurations and checksums. Separate preload throughput
from paced live latency. Use the roadmap's three repetitions for comparable
performance configurations and report mean and standard deviation. Explain
warm-up, drain exclusions, clock uncertainty and failed attempts.}

\subsection{Existing evidence to reuse}
The following retained checks are starting points, not final benchmark claims.
The synthetic late-event check used runtime revision
\texttt{35eb8f6} and verifies injected/exported identities and the matching
oracle exclusions~\cite{late}. It is not full-week or performance evidence.
The full-week gate used runtime revision \texttt{e144981}, with the
\texttt{week-1} and \texttt{week-2} clean replays~\cite{week}. It checks both
queries against the oracle; these are correctness runs, not throughput
measurements. The live latency run at revision \texttt{065eaa6} is
\path{issue37-latency-subset12-live-100x-20261101T210052Z-r1}~\cite{latency}.
It exercises the measurement path on a 12-symbol Monday subset once, without
an oracle comparison. It establishes neither repeatability nor saturation.
Full revisions, configurations and limitations remain in the cited records.

\subsection{Correctness}
\pending{Summarize the retained full-week oracle checks, including exact closing
prices, absolute EMA tolerance, advisory precision/recall, duplicate detection,
late-event identities, all exchanges and weekend coverage. Include sparse and
boundary fixtures, timestamp fallbacks, ties, overnight state and a negative
control. Distinguish binary64/decimal diagnostics from accepted tolerance.}
\begin{table}[ht]
\caption{Planned correctness summary. Every result is pending selection and
verification for the final report.}
\label{tab:correctness}
\begin{tabular}{ll}
\toprule
Check & Final result \\
\midrule
Full-week EMA and advisory comparison & Pending \\
Missing, extra and duplicate outputs & Pending \\
Disorder and beyond-bound late events & Pending \\
Subset publication and final windows & Pending \\
\bottomrule
\end{tabular}
\end{table}

\subsection{Throughput, latency and scalability}
\pending{Add full-week preload throughput and speedup by parallelism; live
latency percentiles under load, with inherent window wait separated from
system delay. Document resource ceilings and the actual bottleneck.}
\pending{Compare exchange count, input rate including peaks, symbol skew,
subscriber count and CPU/TaskManager resources. Report per-subtask imbalance,
CPU, memory, GC, network and lag where observed. A local VM is not evidence of
multi-node or cloud scaling. GPU use is pending applicability, not an implied
measurement.}
\begin{figure}[ht]
\centering
\fbox{\parbox[c][1.15in][c]{0.88\columnwidth}{\centering
Reserved for measured throughput and scalability.\\
Pending figure, raw CSV and run IDs.}}
\caption{Figure slot for comparable repeated runs. No measurements plotted.}
\Description{An empty labeled figure slot reserved for future throughput and scalability measurements.}
\end{figure}
\begin{figure}[ht]
\centering
\fbox{\parbox[c][1.15in][c]{0.88\columnwidth}{\centering
Reserved for latency and resource utilization.\\
Pending figure, clock limits and run IDs.}}
\caption{Figure slot for live latency and resource observations.}
\Description{An empty labeled figure slot reserved for future live latency and resource utilization measurements.}
\end{figure}
\pending{Analyze watermark/late-rate trade-offs and any completed ablations.
List remaining experiments honestly. Preserve the evaluation allocation when
compressing the introduction or implementation text.}

\section{Visualization}
\pending{Select a screenshot from the actual dashboard and explain how it helps
interpret EMA and crossover outputs, skew or replay time. Verify each shown
view against the submitted version. Separate planned views from implemented
ones and operational Grafana metrics from query results.}

\section{Conclusion}
\pending{Answer the motivating question using measured evidence. State remaining
limits, untested recovery/scaling paths and the next experiment. Remove all
pending markers only when their evidence and author text are complete.}
